\documentclass[aspectratio=169,10pt]{beamer}
\usetheme[progressbar=frametitle]{metropolis}
\usepackage{booktabs}
\usepackage{graphicx}

\definecolor{navy}{HTML}{1F3864}
\definecolor{accent}{HTML}{2A78D6}
\definecolor{beh}{HTML}{EB6834}
\definecolor{muted}{HTML}{52514E}
\setbeamercolor{frametitle}{bg=navy,fg=white}
\setbeamercolor{progress bar}{fg=accent,bg=accent!20}
\setbeamercolor{alerted text}{fg=beh}

\setbeamerfont{itemize/enumerate body}{size=\footnotesize}
\setbeamerfont{itemize/enumerate subbody}{size=\scriptsize}
\newcommand{\muted}[1]{{\color{muted}#1}}

\title{Model 2: Gradient Boosting (XGBoost)}
\subtitle{Behavioral Predictors of Weight Change (NHANES 2021--2023)}

\begin{document}

% ---------------------------------------------------------------- Slide 1: methods + performance
\begin{frame}{Model 2: XGBoost --- methods and performance}
\begin{columns}[T,onlytextwidth]
\column{0.45\textwidth}
\small\textbf{Methods}\par\vspace{2pt}
{\footnotesize
\begin{itemize}\setlength\itemsep{1pt}\setlength\parskip{0pt}
  \item \textbf{Model:} XGBoost --- many shallow trees fit in sequence; captures \emph{non-linear} effects and \emph{interactions}
  \item \textbf{Predictor sets:} 15 covariates vs.\ covariates + \alert{12 behaviors}
  \item \textbf{Split:} same 80/20 train/test ($2{,}235 / 559$), stratified on outcome
  \item \textbf{Missing data:} handled natively (no imputation)
  \item \textbf{Tuning:} randomized search, 60 settings $\times$ 5-fold CV $\rightarrow$ depth-3 trees, small learning rate, strong regularization
  \item \textbf{Added value ($\Delta R^2$):} 2,000-sample test bootstrap + repeated 5-fold CV (25 folds)
  \item \textbf{Importance:} permutation (single and grouped), SHAP values and interactions
\end{itemize}}

\column{0.51\textwidth}
\small\textbf{Performance} \muted{\footnotesize (baseline test $R^2 = 0$)}\par\vspace{4pt}
{\footnotesize
\setlength\tabcolsep{5pt}
\begin{tabular}{@{}lccc@{}}
\toprule
Predictors & Test $R^2$ & CV $R^2$ & Test RMSE \\
\midrule
Covariates                     & 0.065 & 0.127 & 8.783 \\
\textbf{+ Behaviors}           & \textbf{0.075} & \textbf{0.138} & \textbf{8.735} \\
\midrule
$\Delta R^2$ (test) [95\% CI]  & \multicolumn{3}{l}{0.010 \;[$-0.015$, $0.033$]} \\
$\Delta R^2$ (CV mean)         & \multicolumn{3}{l}{0.011 \;\muted{(better in 22 of 25 folds)}} \\
\bottomrule
\end{tabular}}
\par\vspace{6pt}
\small\textbf{Grouped permutation importance} \muted{\footnotesize(drop in test $R^2$)}\par\vspace{3pt}
{\footnotesize
\setlength\tabcolsep{5pt}
\begin{tabular}{@{}lc@{}}
BMI 1 yr ago alone            & 0.213 \\
Other 14 covariates together  & 0.043 \\
\alert{All 12 behaviors together} & \alert{0.023} \\
\end{tabular}}
\par\vspace{4pt}
{\scriptsize\muted{Behaviors add a small, fairly consistent gain; the test-set CI includes 0.}}
\end{columns}
\end{frame}

% ---------------------------------------------------------------- Slide 2: what the model learned
\begin{frame}{Model 2: what the model learned}
\begin{columns}[T,onlytextwidth]
\column{0.56\textwidth}
\centering
\includegraphics[width=0.95\linewidth]{m2_slide_shap.pdf}\par
{\scriptsize\muted{SHAP dependence, top four behaviors by mean $|$SHAP$|$ (test set).}\par}\vspace{6pt}
{\scriptsize\muted{\textbf{Caution:} PHQ-9 and sleep have large SHAP values but $\approx$0 test permutation importance, so their patterns may not generalize. All results are associations.}\par}

\column{0.41\textwidth}
\small\textbf{Key findings}\par\vspace{2pt}
{\footnotesize
\begin{itemize}\setlength\itemsep{2pt}\setlength\parskip{0pt}
  \item \textbf{Starting BMI dominates} (mean $|$SHAP$|$ 2.55 vs.\ $\leq$ 0.47 for all others)
  \item \textbf{Behavioral effects are non-linear:}
    \begin{itemize}\scriptsize\setlength\itemsep{0pt}
      \item \alert{Moderate activity:} benefit mainly up to $\sim$150--300 min/wk
      \item \alert{Sedentary time:} threshold at $\sim$5 hr/day
      \item \alert{Short sleep} ($<$6.5 hr): less weight loss
      \item \alert{PHQ-9} $\geq$ 3: less weight loss
    \end{itemize}
  \item \textbf{Most reliable on new data:} moderate activity, then sedentary time, vigorous activity
  \item \textbf{Interactions:} strongest all involve starting BMI (BMI $\times$ moderate activity, BMI $\times$ PHQ-9)
\end{itemize}}
\end{columns}
\end{frame}

\end{document}
